\documentclass[11pt]{article}
\usepackage[margin=1in]{geometry}
\usepackage[T1]{fontenc}
\usepackage[utf8]{inputenc}
\usepackage{hyperref}

\title{Econ 5253 Problem Set 4}
\author{Cheng}
\date{Spring 2026}

\begin{document}
\maketitle

\section*{Data Sources I Would Like to Scrape}
I am most interested in financial and macroeconomic data. One source is the
Federal Reserve Economic Data (FRED) API, which provides time series for
interest rates, inflation, labor markets, and GDP components. A second source
is SEC EDGAR filings (10-K, 10-Q, 8-K), which can be used for event studies,
text analysis of firm disclosures, and accounting-based predictors. A third
source is market microstructure and price data (for example, equities, ETFs, and
options) from API providers that expose historical bars and quote data. These
sources are useful for forecasting, asset pricing, and policy analysis.

\section*{Question 5 (JSON Exercise) Answers}
\begin{enumerate}
  \item I created \texttt{PS4a\_Cheng.R} and used a system call with
  \texttt{wget -O dates.json "URL"} to download the JSON file.

  \item I printed the raw JSON to the console using a system call to
  \texttt{cat dates.json}.

  \item I converted JSON to a list with \texttt{fromJSON("dates.json")} and
  then to a data frame using:
\begin{verbatim}
mydf <- bind_rows(mylist$result[-1])
\end{verbatim}

  \item Object types from my run:
\begin{verbatim}
class(mydf) = c("tbl_df", "tbl", "data.frame")
class(mydf$date) = "character"
\end{verbatim}

  \item The first 10 rows were printed with \texttt{head(mydf, 10)}.
  The resulting table has 6 columns:
  \texttt{date, description, lang, category1, category2, granularity}.
\end{enumerate}

\section*{Question 6 (sparklyr Exercise) Answers}
I created \texttt{PS4b\_Cheng.R} with all requested sparklyr pipeline commands:
\texttt{spark\_connect()}, \texttt{copy\_to()}, \texttt{select()},
\texttt{filter()}, combined select+filter pipeline, \texttt{group\_by()} +
\texttt{summarize()}, and \texttt{arrange()}.

On OSCER, first execution failed due package-version dependency conflicts in the
provided R 4.0.2 environment. The key messages in \texttt{PS4b\_Cheng.Rout}
were:
\begin{verbatim}
namespace 'dplyr' 1.0.9 is being loaded, but >= 1.1.2 is required
there is no package called 'sparklyr'
\end{verbatim}

The assignment note says this part may fail on first try and that students are
not penalized in that case. I therefore kept the full script and submitted it,
along with logs documenting the attempt.

If package compatibility is restored, expected object classes are:
\texttt{class(df1)} as tibble/data.frame and \texttt{class(df)} as a Spark
table object (e.g., \texttt{tbl\_spark}).
\end{document}
